\chapter{Discussion and Future Directions}\label{chapter:discussion}
\thispagestyle{myheadings}

	% ========================================================================
	% GENERAL DISCUSSION
	%
	% This chapter is what your committee reads most closely at the defense.
	% It should synthesise rather than summarise: say what the studies mean
	% together that neither means alone.
	% ========================================================================

	\section{Summary of findings}\label{section:discussion:summary}

		TODO: One short paragraph per results chapter, stating the finding
		rather than restating the approach.

	\section{Interpretation}\label{section:discussion:interpretation}

		TODO: Place the findings in the context of the literature reviewed in
		\cref{chapter:introduction}. Where your results agree with prior work,
		say what they add. Where they disagree, propose why -- differences in
		preparation, species, anesthesia, spatial or temporal resolution -- and
		say which explanation could be tested.

	\section{Limitations}\label{section:discussion:limitations}

		TODO: State the limitations that genuinely constrain the conclusions,
		and for each one say what it does and does not prevent you from
		claiming. A frank, specific limitations section is more persuasive than
		a defensive one, and it is where committees probe hardest.

	\section{Future directions}\label{section:discussion:future}

		TODO: Two or three concrete, feasible experiments -- not a wish list.
		For each, state the question, the approach and the outcome that would
		be informative either way.

	\section{Conclusions}\label{section:discussion:conclusions}

		TODO: Close with a short paragraph on what this dissertation
		establishes.
